%% ma: L11-B15-0047
%% lop: 11
%% chuong: 5
%% bai: 15
%% dang: 11.15.5
%% loai: TN4
%% muc: VDC
%% dap_an: "D"
%% nguon: "tai-lieu-hoc-tap-mon-toan-lop-11-theo-chuong-trinh-giao-duc-pho-thong-2018.pdf (Lớp Toán Thầy Duy, Phạm Lê Duy, THPT Chu Văn An), Chương 3 Bài 1, Phần 1, Câu 36 (tr.327)"
%% kien_thuc: "Lớp 11 Bài 15 (giới hạn của dãy số)"
%% trang_thai: cho-duyet
%% ly_do: "Dạng mới đề xuất (11.15.x), chờ giáo viên duyệt dạng; tài liệu gốc không có lời giải mẫu, đáp án do kho tự giải và kiểm bằng sympy."
%% ghi_chu: "Hai biểu thức liên hợp bậc ba; nhiều bước biến đổi nên xếp VDC."
\begin{ex}
Tính giới hạn $L=\lim\left(\sqrt[3]{8n^3+3n^2-2}+\sqrt[3]{5n^2-8n^3}\right)$.
\choice{$+\infty$}{$-7$}{$\dfrac{53}2$}{\True $\dfrac23$}
\loigiai{\sloppy Viết $\sqrt[3]{8n^3+3n^2-2}-2n=\dfrac{3n^2-2}{A^2+2nA+4n^2}$ với $A=\sqrt[3]{8n^3+3n^2-2}$, giới hạn $=\dfrac{3}{12}=\dfrac14$; và $\sqrt[3]{5n^2-8n^3}+2n=\dfrac{5n^2}{B^2-2nB+4n^2}$ với $B=\sqrt[3]{5n^2-8n^3}\sim-2n$, giới hạn $=\dfrac{5}{12}$. Vậy $L=\dfrac14+\dfrac5{12}=\dfrac23$.}
\end{ex}
